\documentclass[11pt, a4paper]{article}\usepackage[]{graphicx}\usepackage[]{xcolor}
% maxwidth is the original width if it is less than linewidth
% otherwise use linewidth (to make sure the graphics do not exceed the margin)
\makeatletter
\def\maxwidth{ %
  \ifdim\Gin@nat@width>\linewidth
    \linewidth
  \else
    \Gin@nat@width
  \fi
}
\makeatother

\definecolor{fgcolor}{rgb}{0.345, 0.345, 0.345}
\newcommand{\hlnum}[1]{\textcolor[rgb]{0.686,0.059,0.569}{#1}}%
\newcommand{\hlsng}[1]{\textcolor[rgb]{0.192,0.494,0.8}{#1}}%
\newcommand{\hlcom}[1]{\textcolor[rgb]{0.678,0.584,0.686}{\textit{#1}}}%
\newcommand{\hlopt}[1]{\textcolor[rgb]{0,0,0}{#1}}%
\newcommand{\hldef}[1]{\textcolor[rgb]{0.345,0.345,0.345}{#1}}%
\newcommand{\hlkwa}[1]{\textcolor[rgb]{0.161,0.373,0.58}{\textbf{#1}}}%
\newcommand{\hlkwb}[1]{\textcolor[rgb]{0.69,0.353,0.396}{#1}}%
\newcommand{\hlkwc}[1]{\textcolor[rgb]{0.333,0.667,0.333}{#1}}%
\newcommand{\hlkwd}[1]{\textcolor[rgb]{0.737,0.353,0.396}{\textbf{#1}}}%
\let\hlipl\hlkwb

\usepackage{framed}
\makeatletter
\newenvironment{kframe}{%
 \def\at@end@of@kframe{}%
 \ifinner\ifhmode%
  \def\at@end@of@kframe{\end{minipage}}%
  \begin{minipage}{\columnwidth}%
 \fi\fi%
 \def\FrameCommand##1{\hskip\@totalleftmargin \hskip-\fboxsep
 \colorbox{shadecolor}{##1}\hskip-\fboxsep
     % There is no \\@totalrightmargin, so:
     \hskip-\linewidth \hskip-\@totalleftmargin \hskip\columnwidth}%
 \MakeFramed {\advance\hsize-\width
   \@totalleftmargin\z@ \linewidth\hsize
   \@setminipage}}%
 {\par\unskip\endMakeFramed%
 \at@end@of@kframe}
\makeatother

\definecolor{shadecolor}{rgb}{.97, .97, .97}
\definecolor{messagecolor}{rgb}{0, 0, 0}
\definecolor{warningcolor}{rgb}{1, 0, 1}
\definecolor{errorcolor}{rgb}{1, 0, 0}
\newenvironment{knitrout}{}{} % an empty environment to be redefined in TeX

\usepackage{alltt}

\usepackage[top = 0.6 in, bottom = 0.6 in, left = 0.8 in, right = 0.8 in]{geometry}
\usepackage{amsmath, amssymb, amsfonts}

\allowdisplaybreaks[4]

\usepackage{enumerate}
\usepackage{array}
\usepackage{multirow}
\usepackage{dingbat}
\usepackage{fontawesome5}
\usepackage{tasks}
\usepackage{bbding}
\usepackage{undertilde}
\usepackage{twemojis}
% how to use bull's eye ----- \scalebox{2.0}{\twemoji{bullseye}}
\usepackage{fontspec}
\usepackage{fancyhdr}
\usepackage{lastpage}

\pagestyle{fancy}
\fancyhf{}
\fancyfoot[C]{Page \thepage\ of \pageref{LastPage}}
\renewcommand{\headrulewidth}{0pt}

\title{MSMS 408 : Practical 13}
\author{{\fontsize{25}{20} \benyorithafont Ananda Biswas} \\[1em] Exam Roll No. : 24419STC053}
\date{\today}

\newfontface\benyorithafont{Benyoritha-G334D.otf}

\newfontface\myfont{Myfont1-Regular.ttf}[LetterSpace=0.05em]

\newfontface\cbfont{CaveatBrush-Regular.ttf}
% how to use --- \myfont --write text here--

\newfontface\gloriahallelujahfont{GLORIAHALLELUJAH-REGULAR.TTF}
\IfFileExists{upquote.sty}{\usepackage{upquote}}{}
\begin{document}

\maketitle


\section*{\faArrowAltCircleRight[regular] \textcolor{blue}{Question}}

\hspace{1cm} Generate random numbers from $\beta_1(2, 3)$ using $U(0, 1)$ random numbers. For samples of size $20, 50, 80, 100$, 

\begin{enumerate}[(i)]
\item Using the simulated numbers, estimate the mean of the distribution via the standard Monte Carlo method.

\item Improve the estimate by applying the antithetic variables technique.

\item Estimate the variance of both the estimators.

\item Compare the results and comment on the effectiveness of the antithetic variables method in reducing variance.
\end{enumerate}

\section*{\faArrowAltCircleRight[regular] \textcolor{blue}{Theory}}

We have to estimate

$$\theta = E[X] = \dfrac{2}{5} = 0.4, \qquad X \sim \beta_1(2, 3).$$

$\bullet$ To generate random numbers from $\beta_1$ distribution, we recall

$$X \sim U(0, 1) \Rightarrow - \ln X \sim \; \text{Exp}(1).$$

Also, $$\sum \limits_{i = 1}^{n} - \ln X_i \sim \; \text{Gamma}(\text{shape} = n, \text{scale} = 1) \qquad \text{and}$$

$$\dfrac{S_1}{S_1 + S_2} \sim \beta_1(m, n)$$
where $S_1$ and $S_2$ are independent Gamma distributions with shape parameters $m$ and $n$ respectively. \\[0.5em]

$\bullet$ If $x_1, x_2, \ldots, x_n$ is a sample of size $n$ from $X$, Monte Carlo estimate of $\theta = E[X]$ is $$\widehat{\theta} = \dfrac{1}{n} \sum \limits_{i = 1}^{n} x_i.$$

$\bullet$ Estimate of $\theta$ by using antithetic variable will be

$$\widehat{\theta} = \dfrac{1}{n} \sum \limits_{i = 1}^{n/2} [X_i + X_i']$$

where $X_i$ and $X_i'$ are negatively correlated random variables.

\newpage

Here

$$X_i = h(\utilde{u}), \qquad \& \qquad X_i' = h(\mathbf{1} - \utilde{u})$$

with

$$\utilde{u} = (u_1, u_2, \ldots, u_5), \qquad \text{and}$$

$$h(u_1, u_2, \ldots, u_5) = \dfrac{\sum \limits_{i = 1}^{2} - \ln u_i}{\sum \limits_{i = 1}^{2} - \ln u_i + \sum \limits_{i = 3}^{5} - \ln u_i}.$$

\section*{\faArrowAltCircleRight[regular] \textcolor{blue}{R Program}}

\begin{enumerate}

\item[\scalebox{2.0}{\twemoji{keycap: 1}}]





\begin{knitrout}
\definecolor{shadecolor}{rgb}{0.969, 0.969, 0.969}\color{fgcolor}\begin{kframe}
\begin{alltt}
\hldef{h} \hlkwb{<-} \hlkwa{function}\hldef{(}\hlkwc{u}\hldef{,} \hlkwc{m}\hldef{,} \hlkwc{n}\hldef{)\{}
  \hldef{u1} \hlkwb{<-} \hldef{u[}\hlnum{1}\hlopt{:}\hldef{m]}
  \hldef{u2} \hlkwb{<-} \hldef{u[(m} \hlopt{+} \hlnum{1}\hldef{)} \hlopt{:} \hldef{(m} \hlopt{+} \hldef{n)]}

  \hldef{x1} \hlkwb{<-} \hlopt{-}\hlkwd{sum}\hldef{(}\hlkwd{log}\hldef{(u1))}
  \hldef{x2} \hlkwb{<-} \hlopt{-}\hlkwd{sum}\hldef{(}\hlkwd{log}\hldef{(u2))}

  \hldef{x} \hlkwb{<-} \hldef{x1} \hlopt{/} \hldef{(x1} \hlopt{+} \hldef{x2)}
\hldef{\}}
\end{alltt}
\end{kframe}
\end{knitrout}

\begin{knitrout}
\definecolor{shadecolor}{rgb}{0.969, 0.969, 0.969}\color{fgcolor}\begin{kframe}
\begin{alltt}
\hldef{my.rbeta} \hlkwb{<-} \hlkwa{function}\hldef{(}\hlkwc{n}\hldef{,} \hlkwc{shape1}\hldef{,} \hlkwc{shape2}\hldef{,} \hlkwc{uniform.sample}\hldef{)\{}

  \hldef{s} \hlkwb{<-} \hlkwd{c}\hldef{()}

  \hlkwa{for} \hldef{(i} \hlkwa{in} \hlnum{1}\hlopt{:}\hldef{n) \{}
    \hldef{s[i]} \hlkwb{<-} \hlkwd{h}\hldef{(uniform.sample[i,], shape1, shape2)}
  \hldef{\}}

  \hlkwd{return}\hldef{(s)}
\hldef{\}}
\end{alltt}
\end{kframe}
\end{knitrout}

\begin{knitrout}
\definecolor{shadecolor}{rgb}{0.969, 0.969, 0.969}\color{fgcolor}\begin{kframe}
\begin{alltt}
\hldef{sample.sizes} \hlkwb{<-} \hlkwd{c}\hldef{(}\hlnum{20}\hldef{,} \hlnum{50}\hldef{,} \hlnum{80}\hldef{,} \hlnum{100}\hldef{)}
\end{alltt}
\end{kframe}
\end{knitrout}

\begin{knitrout}
\definecolor{shadecolor}{rgb}{0.969, 0.969, 0.969}\color{fgcolor}\begin{kframe}
\begin{alltt}
\hldef{theta.hat.monte.carlo} \hlkwb{<-} \hlkwd{c}\hldef{()}

\hlkwa{for} \hldef{(i} \hlkwa{in} \hlnum{1}\hlopt{:}\hlkwd{length}\hldef{(sample.sizes)) \{}

  \hldef{size} \hlkwb{<-} \hldef{sample.sizes[i]}

  \hldef{u} \hlkwb{<-} \hlkwd{matrix}\hldef{(}\hlkwd{runif}\hldef{(size} \hlopt{*} \hldef{(}\hlnum{2} \hlopt{+} \hlnum{3}\hldef{)),} \hlkwc{nrow} \hldef{= size,} \hlkwc{ncol} \hldef{= (}\hlnum{2} \hlopt{+} \hlnum{3}\hldef{))}

  \hldef{theta.hat.monte.carlo[i]} \hlkwb{<-} \hlkwd{mean}\hldef{(}\hlkwd{my.rbeta}\hldef{(size,} \hlnum{2}\hldef{,} \hlnum{3}\hldef{, u))}
\hldef{\}}
\end{alltt}
\end{kframe}
\end{knitrout}

\begin{knitrout}
\definecolor{shadecolor}{rgb}{0.969, 0.969, 0.969}\color{fgcolor}\begin{kframe}
\begin{alltt}
\hldef{df1} \hlkwb{<-} \hlkwd{data.frame}\hldef{(}\hlkwc{sample.size} \hldef{= sample.sizes,}
                  \hlkwc{theta.hat} \hldef{= theta.hat.monte.carlo)}
\end{alltt}
\end{kframe}
\end{knitrout}




% Table created by stargazer v.5.2.3 by Marek Hlavac, Social Policy Institute. E-mail: marek.hlavac at gmail.com
% Date and time: Tue, Apr 07, 2026 - 22:37:06
\begin{table}[!htbp] \centering 
  \caption{Monte Carlo Estimates} 
  \label{Table 1} 
\begin{tabular}{@{\extracolsep{5pt}} cc} 
\\[-1.8ex]\hline 
\hline \\[-1.8ex] 
sample.size & theta.hat \\ 
\hline \\[-1.8ex] 
$20$ & $0.282318$ \\ 
$50$ & $0.405708$ \\ 
$80$ & $0.384456$ \\ 
$100$ & $0.409716$ \\ 
\hline \\[-1.8ex] 
\end{tabular} 
\end{table} 


\item[\scalebox{2.0}{\twemoji{keycap: 2}}]

\begin{knitrout}
\definecolor{shadecolor}{rgb}{0.969, 0.969, 0.969}\color{fgcolor}\begin{kframe}
\begin{alltt}
\hldef{theta.hat.antithetic} \hlkwb{<-} \hlkwd{c}\hldef{()}

\hlkwa{for} \hldef{(i} \hlkwa{in} \hlnum{1}\hlopt{:}\hlkwd{length}\hldef{(sample.sizes)) \{}
  \hldef{size} \hlkwb{<-} \hldef{sample.sizes[i]}

  \hldef{u} \hlkwb{<-} \hlkwd{matrix}\hldef{(}\hlkwd{runif}\hldef{((size} \hlopt{/} \hlnum{2}\hldef{)} \hlopt{*} \hldef{(}\hlnum{2} \hlopt{+} \hlnum{3}\hldef{)),} \hlkwc{nrow} \hldef{= size} \hlopt{/} \hlnum{2}\hldef{,} \hlkwc{ncol} \hldef{= (}\hlnum{2} \hlopt{+} \hlnum{3}\hldef{))}

  \hldef{x} \hlkwb{<-} \hlkwd{my.rbeta}\hldef{(size} \hlopt{/} \hlnum{2}\hldef{,} \hlnum{2}\hldef{,} \hlnum{3}\hldef{, u)}
  \hldef{x_dash} \hlkwb{<-} \hlkwd{my.rbeta}\hldef{(size} \hlopt{/} \hlnum{2}\hldef{,} \hlnum{2}\hldef{,} \hlnum{3}\hldef{,} \hlnum{1}\hlopt{-}\hldef{u)}

  \hldef{theta.hat.antithetic[i]} \hlkwb{<-} \hlkwd{mean}\hldef{(}\hlkwd{c}\hldef{(x, x_dash))}
\hldef{\}}
\end{alltt}
\end{kframe}
\end{knitrout}

\begin{knitrout}
\definecolor{shadecolor}{rgb}{0.969, 0.969, 0.969}\color{fgcolor}\begin{kframe}
\begin{alltt}
\hldef{df2} \hlkwb{<-} \hlkwd{data.frame}\hldef{(}\hlkwc{sample.size} \hldef{= sample.sizes,}
                  \hlkwc{theta.hat} \hldef{= theta.hat.antithetic)}
\end{alltt}
\end{kframe}
\end{knitrout}


% Table created by stargazer v.5.2.3 by Marek Hlavac, Social Policy Institute. E-mail: marek.hlavac at gmail.com
% Date and time: Tue, Apr 07, 2026 - 22:37:06
\begin{table}[!htbp] \centering 
  \caption{Estimates from Antithetic Variable Technique} 
  \label{Table 2} 
\begin{tabular}{@{\extracolsep{5pt}} cc} 
\\[-1.8ex]\hline 
\hline \\[-1.8ex] 
sample.size & theta.hat \\ 
\hline \\[-1.8ex] 
$20$ & $0.399689$ \\ 
$50$ & $0.391278$ \\ 
$80$ & $0.398930$ \\ 
$100$ & $0.397107$ \\ 
\hline \\[-1.8ex] 
\end{tabular} 
\end{table} 


\newpage

\item[\scalebox{2.0}{\twemoji{keycap: 3}}]

\begin{knitrout}
\definecolor{shadecolor}{rgb}{0.969, 0.969, 0.969}\color{fgcolor}\begin{kframe}
\begin{alltt}
\hldef{B} \hlkwb{<-} \hlnum{1000}
\end{alltt}
\end{kframe}
\end{knitrout}

\begin{knitrout}
\definecolor{shadecolor}{rgb}{0.969, 0.969, 0.969}\color{fgcolor}\begin{kframe}
\begin{alltt}
\hldef{var.monte.carlo} \hlkwb{<-} \hlkwd{c}\hldef{()}

\hlkwa{for} \hldef{(j} \hlkwa{in} \hlnum{1}\hlopt{:}\hlkwd{length}\hldef{(sample.sizes)) \{}

  \hldef{size} \hlkwb{<-} \hldef{sample.sizes[j]}

  \hldef{theta.hat.monte.carlo} \hlkwb{<-} \hlkwd{c}\hldef{()}

  \hlkwa{for} \hldef{(i} \hlkwa{in} \hlnum{1}\hlopt{:}\hldef{B) \{}
    \hldef{u} \hlkwb{<-} \hlkwd{matrix}\hldef{(}\hlkwd{runif}\hldef{(size} \hlopt{*} \hldef{(}\hlnum{2} \hlopt{+} \hlnum{3}\hldef{)),} \hlkwc{nrow} \hldef{= size,} \hlkwc{ncol} \hldef{= (}\hlnum{2} \hlopt{+} \hlnum{3}\hldef{))}

    \hldef{theta.hat.monte.carlo[i]} \hlkwb{<-} \hlkwd{mean}\hldef{(}\hlkwd{my.rbeta}\hldef{(size,} \hlnum{2}\hldef{,} \hlnum{3}\hldef{, u))}
  \hldef{\}}

  \hldef{var.monte.carlo[j]} \hlkwb{<-} \hlkwd{var}\hldef{(theta.hat.monte.carlo)}
\hldef{\}}
\end{alltt}
\end{kframe}
\end{knitrout}

\begin{knitrout}
\definecolor{shadecolor}{rgb}{0.969, 0.969, 0.969}\color{fgcolor}\begin{kframe}
\begin{alltt}
\hldef{df3} \hlkwb{<-} \hlkwd{data.frame}\hldef{(}\hlkwc{sample.size} \hldef{= sample.sizes,}
                  \hlkwc{variance} \hldef{= var.monte.carlo)}
\end{alltt}
\end{kframe}
\end{knitrout}


% Table created by stargazer v.5.2.3 by Marek Hlavac, Social Policy Institute. E-mail: marek.hlavac at gmail.com
% Date and time: Tue, Apr 07, 2026 - 22:37:08
\begin{table}[!htbp] \centering 
  \caption{Variance Monte Carlo Estimates} 
  \label{Table 3} 
\begin{tabular}{@{\extracolsep{5pt}} cc} 
\\[-1.8ex]\hline 
\hline \\[-1.8ex] 
sample.size & variance \\ 
\hline \\[-1.8ex] 
$20$ & $0.001721$ \\ 
$50$ & $0.000795$ \\ 
$80$ & $0.000489$ \\ 
$100$ & $0.000377$ \\ 
\hline \\[-1.8ex] 
\end{tabular} 
\end{table} 


\begin{knitrout}
\definecolor{shadecolor}{rgb}{0.969, 0.969, 0.969}\color{fgcolor}\begin{kframe}
\begin{alltt}
\hldef{var.antithetic} \hlkwb{<-} \hlkwd{c}\hldef{()}

\hlkwa{for} \hldef{(j} \hlkwa{in} \hlnum{1}\hlopt{:}\hlkwd{length}\hldef{(sample.sizes)) \{}

  \hldef{size} \hlkwb{<-} \hldef{sample.sizes[j]}

  \hldef{theta.hat.antithetic} \hlkwb{<-} \hlkwd{c}\hldef{()}

  \hlkwa{for} \hldef{(i} \hlkwa{in} \hlnum{1}\hlopt{:}\hldef{B) \{}
    \hldef{u} \hlkwb{<-} \hlkwd{matrix}\hldef{(}\hlkwd{runif}\hldef{((size} \hlopt{/} \hlnum{2}\hldef{)} \hlopt{*} \hldef{(}\hlnum{2} \hlopt{+} \hlnum{3}\hldef{)),} \hlkwc{nrow} \hldef{= size} \hlopt{/} \hlnum{2}\hldef{,} \hlkwc{ncol} \hldef{= (}\hlnum{2} \hlopt{+} \hlnum{3}\hldef{))}

    \hldef{x} \hlkwb{<-} \hlkwd{my.rbeta}\hldef{(size} \hlopt{/} \hlnum{2}\hldef{,} \hlnum{2}\hldef{,} \hlnum{3}\hldef{, u)}
    \hldef{x_dash} \hlkwb{<-} \hlkwd{my.rbeta}\hldef{(size} \hlopt{/} \hlnum{2}\hldef{,} \hlnum{2}\hldef{,} \hlnum{3}\hldef{,} \hlnum{1}\hlopt{-}\hldef{u)}

    \hldef{theta.hat.antithetic[i]} \hlkwb{<-} \hlkwd{mean}\hldef{(}\hlkwd{c}\hldef{(x, x_dash))}
  \hldef{\}}

  \hldef{var.antithetic[j]} \hlkwb{<-} \hlkwd{var}\hldef{(theta.hat.antithetic)}
\hldef{\}}
\end{alltt}
\end{kframe}
\end{knitrout}

\begin{knitrout}
\definecolor{shadecolor}{rgb}{0.969, 0.969, 0.969}\color{fgcolor}\begin{kframe}
\begin{alltt}
\hldef{df4} \hlkwb{<-} \hlkwd{data.frame}\hldef{(}\hlkwc{sample.size} \hldef{= sample.sizes,}
                  \hlkwc{variance} \hldef{= var.antithetic)}
\end{alltt}
\end{kframe}
\end{knitrout}


% Table created by stargazer v.5.2.3 by Marek Hlavac, Social Policy Institute. E-mail: marek.hlavac at gmail.com
% Date and time: Tue, Apr 07, 2026 - 22:37:10
\begin{table}[!htbp] \centering 
  \caption{Variance Estimates from Antithetic Variable Technique} 
  \label{Table 4} 
\begin{tabular}{@{\extracolsep{5pt}} cc} 
\\[-1.8ex]\hline 
\hline \\[-1.8ex] 
sample.size & variance \\ 
\hline \\[-1.8ex] 
$20$ & $0.000436$ \\ 
$50$ & $0.000173$ \\ 
$80$ & $0.000113$ \\ 
$100$ & $0.000083$ \\ 
\hline \\[-1.8ex] 
\end{tabular} 
\end{table} 


\item[\scalebox{2.0}{\twemoji{keycap: 4}}]

\begin{knitrout}
\definecolor{shadecolor}{rgb}{0.969, 0.969, 0.969}\color{fgcolor}\begin{kframe}
\begin{alltt}
\hldef{relative.gain} \hlkwb{<-} \hldef{((var.monte.carlo} \hlopt{-} \hldef{var.antithetic)} \hlopt{/} \hldef{var.monte.carlo)} \hlopt{*} \hlnum{100}
\end{alltt}
\end{kframe}
\end{knitrout}

\begin{knitrout}
\definecolor{shadecolor}{rgb}{0.969, 0.969, 0.969}\color{fgcolor}\begin{kframe}
\begin{alltt}
\hldef{df5} \hlkwb{<-} \hlkwd{data.frame}\hldef{(}\hlkwc{sample.size} \hldef{= sample.sizes,} \hlkwc{gain} \hldef{= relative.gain)}
\end{alltt}
\end{kframe}
\end{knitrout}


% Table created by stargazer v.5.2.3 by Marek Hlavac, Social Policy Institute. E-mail: marek.hlavac at gmail.com
% Date and time: Tue, Apr 07, 2026 - 22:37:10
\begin{table}[!htbp] \centering 
  \caption{Percentage Variance Reduction} 
  \label{Table 5} 
\begin{tabular}{@{\extracolsep{5pt}} cc} 
\\[-1.8ex]\hline 
\hline \\[-1.8ex] 
sample.size & gain \\ 
\hline \\[-1.8ex] 
$20$ & $74.692$ \\ 
$50$ & $78.175$ \\ 
$80$ & $76.923$ \\ 
$100$ & $77.987$ \\ 
\hline \\[-1.8ex] 
\end{tabular} 
\end{table} 


\smallpencil \hspace{0.1cm} {\fontsize{12}{20}\gloriahallelujahfont We achieve significant reduction in variance of the estimates by using antithetic variable technique.}

\end{enumerate}



\end{document}
